\section{Experimental Evaluation}
\label{sec:eval}
The primary current evidence is R20: six sequential tasks through one
registered harness across two configured simulator scenes, with an
R20-specific source reconstruction kit; it supports bounded task-to-device
integration only. Secondary current evidence is R13: an application-level
two-scene protocol and data-path run whose exact runtime source was not pinned;
it supports protocol, readback, lineage, and gate-refusal behavior. Historical
R1/R2 archives and other protocol, agent, controller, deployment, process,
and recipe records use older versions or direct-node/scripted interfaces and
are contextual software evidence, not pooled process results or matched
strategy comparisons. All process experiments use repository simulators; the
reproducibility notes map each claim to its source and interface version.

\subsection{Evidence Units and Experimental Procedure}
\label{sec:testbed}
\label{sec:procedure}
Historical R1/R2 cases use seed-42 direct-node missions with scripted
proposals; R1's runtime worktree is not archived, so its recorded commit does
not guarantee exact replay. The deterministic controller also uses seeds
43--44. Deployment evidence comprises nine v0.7.57 runs and 782 software
assertions; the recipe case is from the October 18 archive. The same 65
profile/composition assertions were repeated on October 20, 24, and 25 with
mock I/O, not 195 distinct checks or process trials. The cast-film reference
uses an author-defined objective $J$ for thickness conformance, defects,
energy, and throughput.

\subsection{Current-Version Two-Scene Simulator Integration}
\label{sec:crossscene}
An exploratory v0.8.0 application run used the cast-film and annealing
simulator presets. Five protocol families---OPC UA, MQTT, Modbus TCP, Modbus
RTU over TCP, and HTTP/REST---passed read probes; bounded writes returned
matching readbacks. Six 120-s runs per scene at 1-s sampling produced
lineage-verified datasets.

The configured data gate selected the largest absolute signed Pearson
correlation between first-differenced MV and lagged CV over row lags 1--3,
then required $|r_{\Delta}|\geq0.2$, one-step test NRMSE $\leq0.10$, and at
least 20\% four-step rollout NRMSE reduction over persistence. Both
post-pilot candidates failed the gates; neither produced an accepted model,
profile, recommendation, or dispatch. The overlapping 1-s windows are
serially dependent. This run therefore supports protocol/data-path execution
and rejection behavior only, not model quality or process performance.
Signal-pair selection followed an exploratory pilot, and the run-time source
diff was not preserved; the R13 results are included as secondary diagnostic
evidence, with detailed records in the reproducibility supplement.

Table~\ref{tab:r13gates} reports the selected signal pair, the full-dataset
first-difference correlation used for eligibility screening, and two model
metrics computed on held-out test rows. Both scenes failed all three criteria;
these values document the recorded rejection path and are not confirmatory
model-performance estimates.
\begin{table}[!t]
\centering
\caption{R13 exploratory model-gate diagnostics. The correlation column is
the full-dataset eligibility screen; NRMSE and rollout reduction use held-out
test rows. Both candidates fail all three criteria.}
\label{tab:r13gates}
\footnotesize
\setlength{\tabcolsep}{2.2pt}
\begin{tabularx}{\columnwidth}{@{}lXccc@{}}
\toprule
Scene & MV $\to$ CV & $|r_{\Delta}|$ (lag) & 1-step NRMSE & 4-step reduction \\
\midrule
Film & \shortstack[l]{\texttt{ScrewSpeedSP}\\$\to$ \texttt{thick}} & 0.0949 (2) & 0.2587 & 14.11\% \\
Anneal & \shortstack[l]{Z3 setpoint\\$\to$ \texttt{StripTemp}} & 0.0526 (3) & 0.3632 & 4.83\% \\
\bottomrule
\end{tabularx}
\end{table}
\subsection{Registered-Harness Process Path}
\label{sec:harness-e2e}
R20 ran the isolated v0.8.0 working-tree application with the registered
\texttt{omp} adapter and GLM-5.3-Flash in cast-film and annealing simulator
scenes. The app and simulator diffs, base revisions, dependency lockfiles, and
runtime prompt assets are included with the run record; package manifests are
recovered from the recorded base revisions. The external model service does
not provide a deterministic replay of its selection. Each execution created
a separate line, product, recipe, run, channel, agent, and task context; three
sequential executions used one simulator
instance per scene, so they are not independent replicates. Cast film used
OPC UA for screw-speed control and MQTT for the thickness probe; annealing
used Modbus TCP for both the zone-3 setpoint and strip-temperature probe.
The cast-film candidate set was $\{125,130,135\}$~rpm and the annealing
candidate set was $\{735,740,745\}\,^{\circ}$C. Given a live control read,
process measurement, and target band, GLM-5.3-Flash proposed one bounded
\texttt{recipe\_trial}. A synthetic administrator approved each proposal;
the route log recorded one trial and a 1/1 verified device receipt, and all
six tasks reached the completed state. Following task completion and its
evidence queries, a 20-s wait was initiated; one driver probe per task was
then ingested through the isolated SQLite time-series fallback and returned
under the active run tag. The integration checks cover per-execution context, an approved bounded
proposal, route-specific receipt and readback, task completion, and an
active-run-tagged sample; all six records contain this software path. The
runner recorded 76 passing checks and no failures. The task history contains
mirrored event projections, so the per-task operation count is taken from the
route-specific recipe and ACK records.

Table~\ref{tab:q8} reports the pre-task measurement, selected setpoint, and
single post-trial sample. All six post-trial values remained outside the
declared target bands. This experiment verifies the bounded
harness-to-simulator execution and evidence path; it does not show process
attainment, convergence, or quality improvement. The single sample is not a
response trajectory, and the SQLite fallback does not test the normal online
DAQ backend. The two-scene run also does not isolate profile/binding changes
or establish cross-harness portability.

\begin{table}[!t]
\centering
\caption{Model-selected settings and simulator measurements. Each row
reports the pre-task value, configured target band, manipulated variable's
initial-to-selected setting, and one run-tagged sample taken after a 20-s wait
initiated following task completion and evidence queries. All six samples are
outside target; executions are sequential within each scene.}
\label{tab:q8}
\footnotesize
\setlength{\tabcolsep}{2.5pt}
\begin{tabularx}{\columnwidth}{@{}lXcX@{}}
\toprule
Case & Pre-task CV / target & Setting: start$\to$trial & CV after wait \\
\midrule
Film-1 & $55.82/48$--$52~\mu$m & $120\to125$~rpm & $45.69~\mu$m \\
Film-2 & $45.69/48$--$52~\mu$m & $120\to130$~rpm & $46.69~\mu$m \\
Film-3 & $46.69/48$--$52~\mu$m & $120\to125$~rpm & $45.69~\mu$m \\
Anneal-1 & $675.40/700$--$760~^{\circ}$C & $730\to745~^{\circ}$C & $679.60~^{\circ}$C \\
Anneal-2 & $678.30/700$--$760~^{\circ}$C & $730\to745~^{\circ}$C & $677.70~^{\circ}$C \\
Anneal-3 & $676.90/700$--$760~^{\circ}$C & $730\to745~^{\circ}$C & $678.10~^{\circ}$C \\
\bottomrule
\end{tabularx}
\end{table}

\subsection{Heterogeneous Protocol Integration and Route Checks}
\label{sec:protocol}
The suite acquired \resDaqSamples{} samples over five protocol families and
issued policy-checked writes over four; the RTU-over-TCP line was
acquisition-only.
Archived upper-median write latencies ranged from \resPfiftyRange{}, with observed maxima
\resPninetyfiveRange{}. The four writable lines accepted legal writes and
rejected 24/24 out-of-window probes. A configuration-only second preset
added 9/9 rejected probes on three further writable lines. These results
exercise the same resolver against different signal maps.

In a mock-driver API fixture, three deterministic repetitions of six attack
classes per arm isolated admission from endpoint confirmation. The full and
no-readback arms each rejected all 18 probes. The no-interlock and ungated
arms each rejected the 12 hard-range probes but allowed all six soft-window
probes to execute. Thus removing readback did not change admission, while
removing the soft-window gate changed it; confirmation cannot replace a
pre-write gate. A frozen-measurement drill further separated command and
process evidence: the setpoint and readback agreed, yet the monitoring
quantity remained outside its window and triggered backstop rollback after
\resBackstop{}. A manual-binding drill suspended execution until
harness-issued approval (\resHitlLatency{} mediation span).

\subsection{Live-Engine Tasks and Deterministic Controller}
\label{sec:agentteam}
The archived team workflow linked parent task, worker, command, verdict, and
journal. Its single-node cast-film task completed after changing the heater
from \resMissionFrom{} to \resMissionTo\,$^{\circ}$C and observing
\resMissionPv\,$^{\circ}$C within target. Two live-engine tasks used the
same \texttt{omp} harness and GLM-5.3-Flash provider on a scratch line. Their
logs trace a bounded write, descent to \resGoalFinal\,$\mu$m for a
$52.0\pm0.8$\,$\mu$m target, and recipe save; the 52.02\,$\mu$m window
mean differs from the run-summary terminal value. These traces demonstrate
execution integration, not strategy comparison or cross-harness portability.

\label{sec:controller}
A deterministic controller used the dimensionless cast-film objective
\begin{equation}
\begin{aligned}
J={}&55J_{\mathrm{thick}}+25J_{\mathrm{quality}}\\
&+8J_{\mathrm{energy}}+7J_{\mathrm{throughput}},
\end{aligned}
\label{eq:objective}
\end{equation}
whose components score thickness conformance, quality, and energy/throughput
proxies. The latter use screw and line speed, not metered energy or output.
This deterministic controller is not the optional model service.
Table~\ref{tab:q3search} gives its finite grid; formulas and replay details
are in \path{REPRODUCIBILITY.md}.

\begin{table*}[!t]
\centering
\caption{Finite-grid search protocol for the deterministic cast-film
controller. The reference score is restricted to this grid and is not
a global continuous optimum.}
\label{tab:q3search}
\footnotesize
\setlength{\tabcolsep}{4pt}
\begin{tabularx}{\textwidth}{@{}p{2.4cm}Xp{3.0cm}@{}}
\toprule
Search element & Candidate set or rule & Role in the search \\
\midrule
Objective & Eq.~(\ref{eq:objective}); weights 55/25/8/7 for thickness,
quality, energy proxy, and throughput proxy & Rank eligible steady states \\
Eligibility & Melt temperature 195--225~$^{\circ}$C; melt pressure
$\leq 22$~MPa & Discard ineligible states \\
Equal zone setpoint & 195--225~$^{\circ}$C in 2.5~$^{\circ}$C steps & Search variable \\
Screw speed & 50--200~rpm in 5-rpm steps & Search variable \\
Line speed & 20--120~m/min in 2.5-m/min steps & Search variable \\
Die gap & Fixed at 1~mm & Held constant \\
Selection & Evaluate modeled steady state and retain the highest eligible
score; $J^{*}=89.894$ & Finite-grid reference \\
\bottomrule
\end{tabularx}
\end{table*}

Across seeds 42--44, the controller made two policy-checked writes per seed
without rejection and raised $J$ from about \resJzero{} to
$J/J^*=\resControllerMin$--\resControllerMax{} ($J^*=\resJstar$ on the
admissible finite grid). This demonstrates one conventional strategy through
the shared execution services, not comparative control benefit.

\begin{table*}[!t]
\centering
\caption{Recorded injection-recipe actions and post-wait observations
($n=1$ simulator trajectory). Hold time remained 6\,s; each trial used
harness approval, a 65-s wait, and a 1/1 parameter-verification receipt.}
\label{tab:q6actions}
\footnotesize
\setlength{\tabcolsep}{3.5pt}
\begin{tabularx}{\textwidth}{@{}l c c c c X@{}}
\toprule
Stage & Pressure (bar) & Weight (g) & Flash (\%) & Sink (\%) & Recorded action and outcome \\
\midrule
Baseline & 45 & 31.28 & 0 & 1.977 & No command; initial observation \\
Trial 1 & 45$\rightarrow$50 & 31.49 & 0.014 & 1.956 & Approved; 1/1 verified; continue \\
Trial 2 & 50$\rightarrow$55 & 31.74 & 0.031 & 1.760 & Approved; 1/1 verified; continue \\
Trial 3 & 55$\rightarrow$60 & 32.05 & 0 & 1.534 & Weight below 32.15-g minimum and sink above 1.50\%; continue \\
Trial 4 & 60$\rightarrow$65 & 32.30 & 0.019 & 1.277 & Stop recorded; displayed values are within ranges \\
Save & 65 & --- & --- & --- & \texttt{recipe\_update} saved version 2; no \texttt{recipe\_apply} is archived \\
\bottomrule
\end{tabularx}
\end{table*}

\subsection{Deployment Composition and Confirmation Semantics}
\label{sec:hardening}
The archived v0.7.57 deployment suite composed \resLinesLive{} lines,
\resChannelsLive{} channels, \resDaqOnline{} online DAQ nodes,
\resWriteNodesLive{} write nodes, \resRecipesLive{} recipes, and
\resBindingsLive{} bindings on one simulator host. Setpoint/readback/mirror
values agreed within \resConsistencyErr{}, but S4 remained
\texttt{converged=false}; 65-bar observation-to-verification took 96.105\,s.
S5 applied 40-to-40.5$^{\circ}$C, refused a subsequent write, and restored
40$^{\circ}$C after 66.046\,s. Seven injected faults exercised receipt
grading, partial batch failure, and manual confirmation. These software
checks characterize only the tested faults and do not estimate field rates.

\subsection{Historical Four-Process Simulator Missions}
\label{sec:scenarios}
The historical Q4 archive contains one seed-42 scripted, direct-node trajectory
for each of four processes. All four recorded terminal checks met their target
bands, but these runs bypass the current task runtime and do not test
profile/binding adaptation. They are contextual simulator evidence and are
not used to support the current-version integration claims.

\subsection{Operational Recipe Workflow Case Study}
\label{sec:recipeexperiment}
R6 is one October 18 injection-simulator trajectory with four approved
5-bar \texttt{recipe\_trial}s from 45\,bar, each followed by a 65-s wait.
It saved recipe version 2 but did not apply it (Table~\ref{tab:q6actions});
the source commit, seed, and per-sensor timestamps are unavailable, and the
intervention was scripted. R6 bypassed trial cadence and the separate
operation interval. Melt temperature was missing; stopping used weight,
flash, and sink, with freshness based on the newest available timestamp.
Thus the recorded stop is not an independently audited guard pass. A separate
October 23 public-route check refused eight direct-node writes; its
default-cadence recipe run missed target and failed restoration. Neither it
nor the distinct 65-to-70\,bar check replicates R6.
